\subsection{Formula for the pentagon product}
\label{subsec:rotated-regular-polygons-formula-for-the-pentagon-product}

For a real number \(\theta\), define
\[
  d(\theta)
  :=
  \min_{k\in\mathbb{Z}}\left|\theta-\frac{k\pi}{5}\right|.
\]
Then \(0\le d(\theta)\le\pi/10\).

\begin{theorem}[Pentagon product]
\label{thm:rotated-pentagon-product-formula}
Let \(P_5\subset\mathbb{R}^2\) be the regular pentagon with circumradius \(1\),
and let \(K_\theta=P_5\times_L R(\theta)P_5\). Then, for all
\(\theta\in\mathbb{R}\),
\[
  \operatorname{sys}(K_\theta)
  =
  \frac{5+2\sqrt{5}}{10}\sec^2(d(\theta)).
\]
In particular, the maximum in this one-parameter family is attained exactly at
the angles with \(d(\theta)=\pi/10\), where
\[
  \operatorname{sys}(K_\theta)=\frac{3+\sqrt{5}}{5}>1.
\]
These maximizing angles include the Haim--Kislev--Ostrover pentagon product,
up to simultaneous rotations and the rotational symmetries of the pentagon.
\end{theorem}

\begin{proof}
The proof has three parts. First, symmetries reduce the theorem to the
half-domain \(0\le\theta\le\pi/10\). Second, an explicit branch gives the
required upper bound for the capacity, and hence for the systolic ratio. Third,
we reduce the corresponding lower bound on the open half-domain to finitely
many exact algebraic checks, and verify those checks by an executable SageMath
certificate. Continuity fills the remaining cut points and the endpoints. At
the end, we apply the initial symmetry reduction to return from the half-domain
to all real \(\theta\).

The family is periodic with period \(2\pi/5\): simultaneous rotation of both
factors is symplectic, and the two pentagons are invariant under rotation by
\(2\pi/5\). Reflection gives the usual reduction to
\([0,\pi/5]\). There is one further symmetry because the two factors are the
same odd regular polygon: the symplectic map \(-J_0\) swaps the factors and
sends
\[
  P_5\times_L R(\theta)P_5
  \quad\text{to}\quad
  R(\theta)P_5\times_L(-P_5).
\]
After a simultaneous rotation by \(-\theta\), and using
\(-P_5=R(\pi)P_5=R(\pi/5)P_5\) modulo the \(2\pi/5\)-symmetry of the pentagon,
this becomes \(P_5\times_L R(\pi/5-\theta)P_5\). Thus it is enough to prove
\[
  \operatorname{sys}(K_\theta)
  =
  \frac{5+2\sqrt{5}}{10}\sec^2\theta
  \qquad
  \text{for }0\le\theta\le\frac{\pi}{10}.
\]

We now compute the candidate branch on this half-domain. In the raw-sigma
notation used below, this is the branch
\[
  (3,8,1,0,5,6).
\]
With the facet labels introduced below, this is the two-bounce word with
\(q\)-blocks \((3)\) and \((1,0)\), and \(p\)-blocks \((8)\) and \((5,6)\).
Solving the KKT equations for this branch gives the action
\[
  c_{\mathrm{act}}(\theta)
  =
  \left(1+\cos\left(\frac{\pi}{5}\right)\right)^2\sec\theta .
\]
The four-volume is independent of \(\theta\), since rotation preserves area in
the second factor:
\[
  \operatorname{vol}_4(K_\theta)
  =
  \operatorname{area}(P_5)^2,
  \qquad
  \operatorname{area}(P_5)
  =
  \frac{5}{2}\sin\left(\frac{2\pi}{5}\right).
\]
Therefore the active branch gives
\[
  \frac{c_{\mathrm{act}}(\theta)^2}{2\operatorname{vol}_4(K_\theta)}
  =
  \frac{\left(1+\cos\left(\frac{\pi}{5}\right)\right)^4}
       {2\operatorname{area}(P_5)^2}\sec^2\theta
  =
  \frac{5+2\sqrt{5}}{10}\sec^2\theta .
\]
This proves the required upper bound for \(\operatorname{sys}(K_\theta)\) on
the half-domain, since the EHZ capacity is the minimum action.

